\chapter{Number Theory}

\section{Chinese Remaineder Theorem}
\begin{definition}[Chinese Remainder Theorem]
    \label{def:crt}
    co-prime moduli etc.
    \begin{align*}
        x &\equiv x_0 \pmod{q_0} \\
        \vdots \\
        x &\equiv x_{k-1} \pmod{q_{k-1}}
    \end{align*}
    \begin{gather}
        \nonumber
        \text{Let } Q = \prod_{i=0}^{k-1}q_i, \; \hat{q}_i = \frac{Q}{q_i} 
        \\
        \text{Then } x = \sum_{i=0}^{k-1} x_i \cdot \hat{q}_i \cdot \hat{q}_i^{-1} \pmod{Q} 
        \\
        \label{eq:crt-gadget}
        x = \langle \{ x_i \}_{i=0}^{k-1}, \{ \hat{q}_i\cdot\hat{q}_i^{-1} \}_{i=0}^{k-1} \rangle \pmod{Q}
    \end{gather}
\end{definition}

\section{Number Theoretic Transform}
\begin{itemize}
    \item Turns a convolution into a nega/cyclic 
    \item Same as FFT but over different ring
\end{itemize}